\documentclass[12pt,a4paper]{report}

% Bản nháp theo phương án KLTN đại học trong guideline; thay bằng class/template
% chính thức của khoa khi được cung cấp.
\usepackage[a4paper,top=3cm,bottom=3.5cm,left=3.5cm,right=2cm]{geometry}
\usepackage{fontspec}
\IfFontExistsTF{Times New Roman}{\setmainfont{Times New Roman}}{\setmainfont{TeX Gyre Termes}}
\usepackage{polyglossia}
\setdefaultlanguage{vietnamese}
\linespread{1.5}
\AtBeginDocument{\fontsize{13}{19.5}\selectfont}
\setlength{\emergencystretch}{3em}
\usepackage{graphicx}
\usepackage{booktabs}
\usepackage{array}
\usepackage{amsmath}
\usepackage{hyperref}
\hypersetup{colorlinks=true,linkcolor=black,citecolor=black,urlcolor=blue}
\pagestyle{plain}

\newcommand{\todo}[1]{\textbf{[TODO: #1]}}
\newcommand{\thesistitle}{Nhận dạng cử chỉ tay tĩnh ASL sử dụng phân đoạn bàn tay và học chuyển giao}

\begin{document}

\begin{titlepage}
\centering
{\large ĐẠI HỌC QUỐC GIA THÀNH PHỐ HỒ CHÍ MINH\\
\textbf{TRƯỜNG ĐẠI HỌC CÔNG NGHỆ THÔNG TIN}}\\[2cm]
{\Large \textbf{ĐỒ ÁN TỐT NGHIỆP}}\\[1.5cm]
{\LARGE \textbf{\thesistitle}}\\[2cm]
\begin{flushleft}
\large
Sinh viên thực hiện: Nguyễn Hoàng Nam\\
\hspace*{4.25cm}Trần Ngọc Sơn\\[0.5cm]
Giảng viên hướng dẫn: \todo{Bổ sung tên giảng viên}
\end{flushleft}
\vfill
{\large Thành phố Hồ Chí Minh, 2026}
\end{titlepage}

\chapter*{LỜI CAM ĐOAN}
\addcontentsline{toc}{chapter}{LỜI CAM ĐOAN}
Nhóm chúng tôi cam đoan nội dung đồ án này là kết quả nghiên cứu và triển khai của nhóm dưới sự hướng dẫn của giảng viên. Các kết quả, dữ liệu và tài liệu kế thừa từ nguồn khác đều được ghi rõ nguồn trích dẫn.\todo{Rà soát và ký theo mẫu chính thức.}

\chapter*{TÓM TẮT}
\addcontentsline{toc}{chapter}{TÓM TẮT}
Đồ án trình bày các baseline tái lập cho bài toán phân loại 36 lớp ảnh bàn tay theo quy ước của tập dữ liệu ASL-HG, gồm A--Z và 0--9. Protocol hiện hành dùng archive ảnh đã xử lý, loại trùng lặp chính xác và chia dữ liệu theo người ký hiệu. Mô hình tốt nhất hiện có là MobileNetV4 Conv-S tiền huấn luyện trên ImageNet và tinh chỉnh toàn bộ tầng; các số liệu chỉ được diễn giải trong giới hạn split đã khóa. MediaPipe Hand Landmarker được đánh giá riêng ở baseline landmarks+SVM; hướng định vị ROI bằng MediaPipe kết hợp MobileNetV4 là công việc tiếp theo, không được báo cáo như kết quả đã hoàn thành. Với webcam, mọi mô hình tương lai vẫn chỉ phân loại từng khung hình độc lập và chưa mô hình hóa thông tin thời gian.\todo{Bổ sung benchmark realtime khi có artifact tương ứng.}

\chapter*{DANH MỤC TỪ VIẾT TẮT}
\addcontentsline{toc}{chapter}{DANH MỤC TỪ VIẾT TẮT}
\begin{tabular}{>{\bfseries}p{3cm}p{10cm}}
ASL & American Sign Language\\
CNN & Convolutional Neural Network\\
GPU & Graphics Processing Unit\\
FPS & Frames Per Second\\
\end{tabular}

\tableofcontents
\listoftables
\listoffigures

\chapter*{MỞ ĐẦU}
\addcontentsline{toc}{chapter}{MỞ ĐẦU}

\section*{1. Lý do chọn đề tài}
\addcontentsline{toc}{section}{1. Lý do chọn đề tài}
American Sign Language (ASL) là một ngôn ngữ ký hiệu được sử dụng trong cộng đồng người khiếm thính. Nhận dạng lớp ảnh bàn tay là một bài toán thị giác máy tính có thể phục vụ nghiên cứu, học tập về phân loại ảnh và thử nghiệm tương tác người--máy. Trong phạm vi đồ án, mỗi ảnh hoặc khung hình đầu vào được gán một nhãn riêng theo quy ước của tập dữ liệu; hệ thống không được xem là hệ thống nhận dạng ngôn ngữ ký hiệu hoàn chỉnh hoặc hệ thống dịch giao tiếp.

Hiệu quả của mô hình phụ thuộc vào chất lượng dữ liệu, vị trí bàn tay, điều kiện ánh sáng và mức độ tương đồng về hình dạng giữa các lớp. Các baseline ảnh hiện tại dùng archive ảnh đã xử lý của ASL-HG; chúng không chạy MediaPipe trong khi huấn luyện. MediaPipe Hand Landmarker được đánh giá riêng ở baseline landmarks+SVM. Một pipeline ROI dựa trên landmark, nếu được mở rộng, chỉ là định vị và trích xuất vùng quan tâm (region of interest -- ROI), không phải phân đoạn theo từng điểm ảnh. Sự nhầm lẫn giữa hai lớp ảnh mang nhãn O và 0 được phân tích riêng bằng dữ liệu đánh giá, không được giả định là lỗi phổ biến của mọi hệ thống trước đó.

\section*{2. Mục tiêu nghiên cứu}
\addcontentsline{toc}{section}{2. Mục tiêu nghiên cứu}
Mục tiêu tổng quát của đồ án là xây dựng và đánh giá pipeline phân loại 36 lớp ảnh bàn tay theo cách biểu diễn của tập ASL-HG. Các mục tiêu cụ thể gồm:
\begin{itemize}
    \item xây dựng quy trình tái lập được từ dữ liệu đã khóa phiên bản, kiểm kê dữ liệu, loại trùng lặp chính xác, chia dữ liệu theo người ký hiệu, huấn luyện và đánh giá;
    \item đánh giá các baseline CNN từ đầu, MediaPipe landmarks kết hợp SVM, và MobileNetV4 Conv-S tiền huấn luyện trên ImageNet ở hai chế độ backbone đóng băng và tinh chỉnh toàn bộ;
    \item đánh giá mô hình bằng Accuracy, Precision, Recall, F1-score theo lớp, confusion matrix và các chỉ số nhầm lẫn theo hai hướng O$\rightarrow$0, 0$\rightarrow$O;
    \item lựa chọn mô hình dựa trên validation set; chỉ sử dụng test set đã khóa cho đánh giá cuối cùng;
    \item xây dựng và đánh giá pipeline MediaPipe-guided ROI kết hợp MobileNetV4.
\end{itemize}

\section*{3. Đối tượng và phạm vi nghiên cứu}
\addcontentsline{toc}{section}{3. Đối tượng và phạm vi nghiên cứu}
Đối tượng nghiên cứu là 36.000 ảnh thuộc 36 lớp A--Z và 0--9 của tập ASL-HG \cite{asl-hg-original}. Repository Hugging Face của nhóm là nguồn phân phối dùng trong quá trình tái lập \cite{asl-hg-repo}. Cụm ``36 lớp ảnh'' phản ánh quy ước nhãn của dataset, không khẳng định toàn bộ các chữ cái ASL đều là cử chỉ tĩnh về mặt ngôn ngữ học: fingerspelling của J và Z có thành phần chuyển động.

Đề tài xét ảnh tĩnh một bàn tay. Nhận dạng cử chỉ động, chuỗi ký hiệu, câu và diễn giải ngữ nghĩa nằm ngoài phạm vi nghiên cứu.

\section*{4. Phương pháp nghiên cứu tổng quát}
\addcontentsline{toc}{section}{4. Phương pháp nghiên cứu tổng quát}
Đồ án áp dụng chiến lược ``Baseline First, Improve Later''. Trước hết, nhóm xây dựng các baseline tái lập từ archive ảnh đã xử lý của dataset đến báo cáo đánh giá. Sau đó, confusion matrix, các chỉ số theo lớp và log thực nghiệm được dùng để lựa chọn cải tiến về dữ liệu, mô hình và tốc độ suy luận. Các notebook được thực thi trên Google Colab thông qua Colab CLI; cấu hình, seed, checkpoint và artifact được lưu để hỗ trợ tái lập.

\section*{5. Đóng góp}
\addcontentsline{toc}{section}{5. Đóng góp}
Các đóng góp có thể kiểm chứng của đồ án là: protocol participant-disjoint có loại trùng lặp SHA-256; các baseline tái lập; pipeline raw image--MediaPipe ROI--MobileNetV4 full fine-tuning; và phân tích O/0. Run 	exttt{mp-mnv4-003} có artifact bất biến, đạt Accuracy 96,38\% và Macro-F1 95,64\% trên P9. Nguyên mẫu desktop chưa là đóng góp hoàn tất vì chưa benchmark.

\section*{6. Bố cục đồ án}
\addcontentsline{toc}{section}{6. Bố cục đồ án}
Chương 1 trình bày tổng quan bài toán và các nghiên cứu liên quan. Chương 2 giới thiệu cơ sở lý thuyết, dữ liệu và các chỉ số đánh giá. Chương 3 mô tả pipeline, mô hình và thiết kế thực nghiệm. Chương 4 trình bày kết quả và bàn luận. Chương 5 tổng kết, nêu hạn chế và hướng phát triển.

\chapter{TỔNG QUAN}

\section{Bài toán nhận dạng ASL}
Nhận dạng ASL từ ảnh trong đồ án này là bài toán ánh xạ một ảnh đầu vào $x$ vào một nhãn $y$ trong tập 36 lớp của ASL-HG. Nhãn biểu diễn chữ cái hoặc chữ số theo quy ước dữ liệu; không nên diễn giải tập nhãn này như khẳng định rằng mọi chữ cái fingerspelling ASL đều tĩnh. Đặc biệt, J và Z có thành phần chuyển động trong fingerspelling ASL. Bài toán không bao gồm nhận dạng động tác theo thời gian hoặc ghép các ký hiệu thành câu.

\section{Tổng quan phân loại ảnh}
Phân loại ảnh là nhiệm vụ dự đoán lớp của ảnh dựa trên các đặc trưng thị giác. Các hệ thống học sâu thường học đồng thời phép biểu diễn đặc trưng và bộ phân loại. Với dữ liệu có kích thước vừa phải, học chuyển giao cho phép sử dụng backbone đã học từ ImageNet rồi điều chỉnh lớp phân loại cho tập lớp đích.

\section{Các hướng tiếp cận liên quan}
Các hướng tiếp cận chính cho nhận dạng cử chỉ gồm sử dụng đặc trưng hình ảnh trực tiếp, sử dụng điểm mốc bàn tay, hoặc kết hợp định vị vùng bàn tay với mạng nơ-ron tích chập. MediaPipe Hand Landmarker sử dụng bộ phát hiện lòng bàn tay để xác định vùng bàn tay và mô hình landmark để suy ra 21 điểm mốc; các điểm này có thể được dùng làm đặc trưng hoặc để định vị ROI \cite{mediapipe}. Trong các thí nghiệm hiện có, nhóm đánh giá CNN từ đầu, MediaPipe landmarks kết hợp SVM, và MobileNetV4 Conv-S tiền huấn luyện \cite{mobilenetv4}.

\section{Vấn đề còn tồn tại}
Hiệu quả của hệ thống có thể suy giảm khi ảnh bị mờ, nền phức tạp hoặc bộ phát hiện không tìm thấy bàn tay. Ngoài ra, hai lớp O và 0 có thể có hình dạng gần nhau trong một số ảnh, khiến Accuracy tổng thể không phản ánh đầy đủ chất lượng từng lớp. Vì vậy, đồ án cần ghi nhận các trường hợp không định vị được bàn tay thay vì âm thầm loại bỏ và cần phân tích các ô nhầm lẫn O--0 riêng biệt.

\section{Khoảng trống và hướng giải quyết}
Đồ án tập trung vào một quy trình có thể tái lập: khóa revision dữ liệu, kiểm kê dữ liệu, canonicalization theo SHA-256, chia participant-disjoint và đánh giá theo cùng protocol. Cách tổ chức này tạo cơ sở để so sánh các baseline và trạng thái backbone đóng băng với tinh chỉnh toàn bộ một cách công bằng. Hướng định vị ROI bằng MediaPipe sẽ chỉ được so sánh với baseline ảnh sau khi policy xử lý trường hợp không phát hiện bàn tay được khóa trước khi xem test metrics.

\section{Kết luận chương}
Chương này đã xác định phạm vi bài toán, các hướng tiếp cận chính và các rủi ro cần đánh giá. Chương tiếp theo trình bày cơ sở lý thuyết và dữ liệu được sử dụng.

\chapter{CƠ SỞ LÝ THUYẾT VÀ DỮ LIỆU}

\section{Biểu diễn bài toán phân loại ảnh}
Với tập dữ liệu $\mathcal{D}=\{(x_i,y_i)\}_{i=1}^{N}$, mô hình nhận một ảnh $x_i$ và sinh ra vector xác suất trên 36 lớp. Nhãn dự đoán được chọn theo xác suất lớn nhất. Trong quá trình huấn luyện, hàm mất mát cross-entropy được tối thiểu hóa trên tập huấn luyện.

\section{Mạng nơ-ron tích chập}
Mạng nơ-ron tích chập (CNN) sử dụng các phép tích chập để học những đặc trưng từ cục bộ đến trừu tượng. Các tầng sâu hơn thường biểu diễn hình dạng và cấu trúc có ích cho phân loại. Trong đồ án, CNN được xem là nền tảng lý thuyết cho backbone học chuyển giao.

\section{Học chuyển giao và MobileNetV4}
Học chuyển giao sử dụng trọng số được huấn luyện trước trên một tập dữ liệu lớn, sau đó thay thế đầu phân loại để phù hợp với tập lớp đích. Mục đích là tận dụng đặc trưng tiền huấn luyện và giảm chi phí huấn luyện so với huấn luyện từ đầu; phương pháp này không tự bảo đảm một mức hiệu năng cụ thể.

Các thí nghiệm ảnh sử dụng MobileNetV4 Conv-S tiền huấn luyện trên ImageNet-1K \cite{mobilenetv4}. Run \texttt{mnv4-001} đóng băng toàn bộ mạng nền và chỉ huấn luyện đầu phân loại 36 lớp. Run \texttt{mnv4-002} khởi tạo từ cùng trọng số ImageNet, thay đầu phân loại rồi mở khóa toàn bộ tầng để tinh chỉnh với learning rate $10^{-4}$. Việc thứ hai là full fine-tuning; nó khác với huấn luyện riêng đầu phân loại ở run đóng băng.

\section{Dữ liệu ASL-HG}
Tập dữ liệu được sử dụng là ASL-HG, gồm 36.000 ảnh thuộc 36 lớp A--Z và 0--9 \cite{asl-hg-original}. Nguồn tải trực tiếp phục vụ tái lập là repository Hugging Face của nhóm \cite{asl-hg-repo}. Đây là quy ước phân lớp của dataset; mô hình không xác nhận tính tĩnh của mọi chữ cái ASL về mặt ngôn ngữ học. Kiểm kê dữ liệu đọc được 36.000 ảnh và phát hiện 994 dòng thuộc 435 nhóm ảnh trùng chính xác. Sau canonicalization theo SHA-256, 35.441 ảnh được giữ lại cho protocol thực nghiệm.

\begin{table}[ht]
\centering
\caption{Thông tin tổng quan về tập dữ liệu ASL-HG}
\label{tab:dataset-overview}
\begin{tabular}{ll}
\toprule
Thuộc tính & Giá trị\\
\midrule
Số lượng ảnh & 36.000\\
Số lớp & 36\\
Nhãn & A--Z và 0--9 theo quy ước ASL-HG\\
Kiểu dữ liệu & Ảnh một bàn tay; phân loại từng ảnh\\
Nguồn tải & Hugging Face Dataset repository\\
\bottomrule
\end{tabular}
\end{table}

\section{MediaPipe Hand Landmarker và tiền xử lý}
MediaPipe Hand Landmarker phát hiện bàn tay và suy ra 21 điểm mốc trên bàn tay \cite{mediapipe}. Trong run \texttt{mp-svm-001}, đặc trưng landmark tối đa hai tay có 130 chiều được chuẩn hóa bằng StandardScaler và phân loại bằng RBF SVM. Với MobileNetV4, các run hiện có dùng archive ảnh đã xử lý của dataset (\texttt{ASL\_Processed\_Images.zip}), đổi kích thước theo data configuration của mô hình và chuẩn hóa bằng mean/std ImageNet; chúng không chạy lại MediaPipe khi huấn luyện.

Run 	exttt{mp-mnv4-003} dùng raw image canonical, MediaPipe Hand Landmarker một tay, ROI vuông padding 0,18 và raw-image fallback khi detection/ROI thất bại. Nó đạt 96,66\% Accuracy và 95,75\% Macro-F1 trên P9, cao hơn 	exttt{mnv4-002} 2,15 điểm Accuracy. Đây là phát hiện/định vị và trích xuất ROI, không gọi là phân đoạn theo từng điểm ảnh vì không có segmentation mask.

\section{Chia dữ liệu}
Dữ liệu được chia theo người ký hiệu (participant-disjoint), không phải random split từng ảnh. Sau canonicalization, tập huấn luyện gồm P1, P3--P8 và P10 (28.365 ảnh); tập xác thực là P2 (3.487 ảnh); tập kiểm thử là P9 (3.589 ảnh). Seed được cố định là 42. Tập xác thực được dùng cho lựa chọn checkpoint và siêu tham số. Tập kiểm thử được giữ nguyên trong quá trình phát triển và chỉ dùng sau khi model selection hoàn tất. Protocol này ngăn cùng SHA-256 xuất hiện ở nhiều split, nhưng chỉ có một participant ở validation/test nên chưa thay thế được leave-one-participant-out hoặc external test set.

\section{Chỉ số đánh giá}
Các chỉ số gồm Accuracy, Precision, Recall và F1-score theo lớp. Confusion matrix được sử dụng để quan sát các cặp lớp dễ nhầm, đặc biệt là hai hướng O$\rightarrow$0 và 0$\rightarrow$O.

\section{Kết luận chương}
Chương này đã trình bày nền tảng CNN, học chuyển giao, dữ liệu và protocol tiền xử lý. Chương 3 mô tả cách các thành phần được kết hợp thành pipeline thực nghiệm.

\chapter{PHƯƠNG PHÁP VÀ THIẾT KẾ THỰC NGHIỆM}

\section{Kiến trúc tổng thể của pipeline}
Quy trình thực nghiệm hiện có gồm: tải archive đã khóa revision, kiểm kê và canonicalization theo SHA-256, tái dựng participant split, tiền xử lý theo data configuration, huấn luyện, lựa chọn checkpoint theo validation và đánh giá test. Notebook là đơn vị thực nghiệm tái lập được.

\begin{figure}[ht]
\centering
\begin{tabular}{c}
\fbox{Archive đã khóa revision} $\rightarrow$ \fbox{Kiểm kê và canonicalization}\\[0.4cm]
$\rightarrow$ \fbox{Participant split} $\rightarrow$ \fbox{Tiền xử lý MobileNetV4}\\[0.4cm]
$\rightarrow$ \fbox{Chọn checkpoint theo validation} $\rightarrow$ \fbox{Đánh giá test}
\end{tabular}
\caption{Quy trình xử lý và đánh giá của hệ thống}
\label{fig:pipeline}
\par\smallskip\textit{Nguồn: nhóm tác giả.}
\end{figure}

\section{Chuẩn bị dữ liệu}
Kiểm kê kiểm tra số lớp, khả năng đọc ảnh, hash và ảnh trùng. Các run dùng một canonical representative cho mỗi SHA-256 để loại leakage. Run \texttt{mp-mnv4-003} lấy raw image canonical theo manifest/split khóa, dùng MediaPipe Hand Landmarker một tay tạo ROI vuông padding 0,18, rồi dùng raw image fallback nếu detection/ROI thất bại. Crop manifest, coverage, checkpoint và metrics được publish cùng artifact.

\section{Mô hình baseline}
Run \texttt{mnv4-001} sử dụng MobileNetV4 Conv-S tiền huấn luyện trên ImageNet với toàn bộ mạng nền bị đóng băng; chỉ đầu phân loại được thay để dự đoán 36 lớp. Run \texttt{mnv4-002} khởi tạo lại từ checkpoint ImageNet bất biến, thay đầu phân loại và cho phép toàn bộ tham số được cập nhật. Cả hai run đều chọn checkpoint có validation loss tốt nhất; do đó, kết quả test không được dùng để quyết định epoch hoặc siêu tham số.

\section{Cấu hình huấn luyện}
\begin{table}[ht]
\centering
\caption{Cấu hình huấn luyện của run \texttt{mp-mnv4-003}}
\label{tab:training-config}
\begin{tabular}{ll}
\toprule
Thành phần & Cấu hình\\
\midrule
Kích thước ảnh & 224$\times$224\\
Backbone & MobileNetV4 Conv-S, ImageNet-1K pretrained\\
Số lớp đầu ra & 36\\
Optimizer & AdamW\\
Learning rate (full fine-tuning) & $10^{-4}$\\
Batch size & 64\\
Số epoch tối đa & 15\\
Seed & 42\\
\bottomrule
\end{tabular}
\end{table}

\section{Thiết kế cải tiến}
CNN, MobileNetV4 frozen, MediaPipe landmarks+SVM, \texttt{mnv4-002} và \texttt{mp-mnv4-003} dùng cùng canonicalization và participant split. \texttt{mp-mnv4-003} fine-tune toàn bộ backbone với learning rate $10^{-4}$; policy fallback được khóa trước khi đọc test. So sánh \texttt{mnv4-002} (archive processed) với \texttt{mp-mnv4-003} (raw image--MediaPipe ROI) là ablation đầu vào hiện có; one-hand/two-hand và external evaluation còn là hướng mở rộng.

\section{Bảo đảm khả năng tái lập}
Mỗi run cần lưu seed, config, phiên bản thư viện, checkpoint, learning curves, confusion matrix và failure manifest. Colab CLI được dùng để cấp GPU và thực thi notebook.

\section{Kết luận chương}
Chương này đã mô tả pipeline baseline, cấu hình mô hình và nguyên tắc thiết kế thí nghiệm. Các kết quả định lượng sẽ được trình bày ở chương tiếp theo sau khi hoàn thành các run chính thức.

\chapter{KẾT QUẢ VÀ BÀN LUẬN}

\section{Môi trường thực nghiệm}
Thí nghiệm được thực thi trên Google Colab thông qua Colab CLI. Mã nguồn notebook, cấu hình và artifact được quản lý trong repository của dự án.\todo{Bổ sung loại GPU, phiên bản TensorFlow/MediaPipe, thời điểm run và experiment ID.}

\section{Kết quả baseline}
Sau khi checkpoint và siêu tham số được chọn bằng validation P2, các baseline được đánh giá trên test P9. Bảng \ref{tab:baseline-results} tổng hợp các artifact đã tái lập. Các con số chỉ áp dụng cho protocol participant-disjoint đã mô tả ở Chương 2.

\begin{table}[ht]
\centering
\small
\caption{Kết quả baseline trên tập kiểm thử}
\label{tab:baseline-results}
\begin{tabular}{p{4.6cm}rrrr}
\toprule
Mô hình & Accuracy & Precision & Recall & F1-score\\
\midrule
CNN từ đầu (\texttt{cnn-001}) & 83,25\% & 79,57\% & 83,31\% & 80,10\%\\
MobileNetV4 Conv-S đóng băng (\texttt{mnv4-001}) & 79,88\% & 77,65\% & 79,94\% & 77,29\%\\
MediaPipe landmarks + RBF SVM (\texttt{mp-svm-001}) & 90,53\% & 87,04\% & 90,52\% & 88,03\%\\
MobileNetV4 Conv-S full fine-tuning (\texttt{mnv4-002}) & \textbf{94,51\%} & \textbf{94,05\%} & \textbf{94,53\%} & \textbf{93,38\%}\\
\bottomrule
\end{tabular}
\end{table}

\section{So sánh các mô hình}
MobileNetV4 full fine-tuning là baseline ảnh mạnh nhất hiện có. So với CNN từ đầu, test accuracy tăng 11,26 điểm phần trăm; so với MediaPipe landmarks+SVM, test accuracy tăng 3,98 điểm phần trăm. MobileNetV4 đóng băng toàn bộ mạng nền không vượt CNN từ đầu trong protocol này. Kết quả cho thấy lợi ích của trọng số ImageNet phụ thuộc vào chiến lược tinh chỉnh, không cho phép kết luận về mọi dataset hoặc mọi phương án transfer learning.

\section{Phân tích confusion matrix và cặp O/0}
Confusion matrix được dùng để xác định các cặp lớp có lỗi cao. Các baseline \texttt{mnv4-001}, \texttt{mp-svm-001} và \texttt{mnv4-002} đều có Recall 100\% ở hai lớp O và 0, không có nhầm lẫn O$\rightarrow$0 hoặc 0$\rightarrow$O trên test P9. Ngược lại, \texttt{cnn-001} có Recall lớp 0 bằng 0\%. Phân tích này chỉ là đánh giá trên split P9, không được dùng để khái quát về mọi hệ thống ASL.\todo{Chèn confusion matrix từ artifact \texttt{mnv4-002}.}

\section{Phân tích lỗi}
Các lỗi cần được phân nhóm theo nguyên nhân quan sát được: không phát hiện bàn tay, ảnh mờ hoặc chất lượng kém, nền/ánh sáng, ROI không phù hợp và hình dạng bàn tay tương tự. Không được suy luận nguyên nhân chỉ từ một ảnh nếu chưa kiểm tra mẫu lỗi.

\section{Bàn luận}
Kết quả hiện có trả lời các mục tiêu về phân loại ảnh, confusion matrix và O/0. Chúng chưa trả lời mục tiêu triển khai thời gian thực vì chưa có benchmark latency/FPS trên phần cứng xác định. Validation và test chỉ gồm một participant mỗi split, nên chênh lệch giữa chúng có thể phản ánh khác biệt người ký hiệu; kết quả không thay thế leave-one-participant-out hoặc external evaluation. Kết quả quan sát được phải được tách khỏi diễn giải của nhóm.

\section{Kết luận chương}
Chương này đã báo cáo các baseline có artifact tái lập. Các thí nghiệm MediaPipe-guided crop kết hợp MobileNetV4, benchmark webcam và đánh giá tổng quát hóa rộng hơn vẫn là công việc tiếp theo.

\chapter{HỆ THỐNG ỨNG DỤNG VÀ TRIỂN KHAI}

\section{Yêu cầu hệ thống}
Nguyên mẫu desktop Python/OpenCV là hạng mục triển khai tiếp theo. Khi được hoàn thành, hệ thống sẽ nhận từng khung hình từ webcam, dùng MediaPipe Hand Landmarker để định vị bàn tay và trích xuất ROI, chạy mô hình phân loại, rồi hiển thị nhãn dự đoán cùng trạng thái phát hiện. Khi không phát hiện được bàn tay hoặc ROI không hợp lệ, hệ thống cần thông báo trạng thái thay vì trả về nhãn dự đoán. Ứng dụng không thực hiện nhận dạng ngôn ngữ ký hiệu hoàn chỉnh hay phân tích chuỗi ký hiệu.

\section{Pipeline suy luận}
Pipeline thời gian thực dự kiến tái sử dụng các bước tiền xử lý của mô hình đã chọn. Mô hình chỉ được nạp một lần khi khởi động; từng khung hình được xử lý độc lập. Thời gian cho MediaPipe Hand Landmarker và thời gian suy luận mô hình cần được đo riêng để xác định nút thắt hiệu năng. Hệ thống chưa dùng mô hình theo thời gian, tracking liên khung hình hoặc cơ chế làm mượt nhãn.

\section{Đánh giá thời gian thực}
Các đại lượng cần ghi nhận gồm độ trễ trung bình, percentile độ trễ, FPS và mức sử dụng tài nguyên. Các kết quả được báo cáo cùng cấu hình thiết bị, độ phân giải khung hình và số khung hình đo.\todo{Bổ sung kết quả benchmark trên thiết bị chạy demo.}

\section{Hạn chế triển khai}
Điều kiện ánh sáng, góc bàn tay, nền ảnh và khả năng phát hiện của MediaPipe có thể làm giảm chất lượng hệ thống. Các hướng phát triển gồm giảm kích thước khung hình, tracking giữa các khung hình, TFLite hoặc quantization, và mô hình hóa chuỗi thời gian nếu mở rộng sang cử chỉ động.

\section{Kết luận chương}
Chương này mô tả nguyên mẫu realtime và các tiêu chí cần kiểm tra. Đây là phần gắn kết kết quả mô hình với ứng dụng thực tế.

\chapter*{KẾT LUẬN VÀ KIẾN NGHỊ}
\addcontentsline{toc}{chapter}{KẾT LUẬN VÀ KIẾN NGHỊ}

\section*{1. Kết luận}
Đồ án xây dựng protocol tái lập cho bài toán phân loại 36 lớp ảnh theo quy ước ASL-HG. Các baseline hiện có sử dụng canonicalization theo SHA-256, participant-disjoint split và MobileNetV4 Conv-S full fine-tuning; hướng định vị/trích xuất ROI bàn tay dựa trên MediaPipe kết hợp mô hình ảnh là công việc tiếp theo.\todo{Cập nhật kết luận khi hoàn thành các thí nghiệm đề xuất.}

\section*{2. Đóng góp}
Các đóng góp đã có gồm artifact tái lập với dataset revision, checksum, canonical split, checkpoint và metrics của các baseline; phân tích riêng O/0; và MobileNetV4 Conv-S full fine-tuning là baseline ảnh tốt nhất trên split P9. Các đóng góp liên quan đến MediaPipe-guided crop và nguyên mẫu desktop chỉ được giữ lại sau khi có artifact tương ứng.

\section*{3. Hạn chế}
Các hạn chế cần đánh giá gồm phạm vi phân loại ảnh hoặc từng khung hình độc lập, không mô hình hóa chuyển động của J/Z hay dynamic gestures, phụ thuộc vào MediaPipe, điều kiện dataset cố định và khả năng tổng quát sang môi trường thực tế khác.\todo{Bổ sung hạn chế dựa trên error analysis.}

\section*{4. Hướng phát triển}
Các hướng phát triển gồm fine-tuning có kiểm soát, augmentation chuyên biệt cho O/0, xử lý nhiều điều kiện ánh sáng, tối ưu mô hình cho thiết bị biên và mở rộng từ nhận dạng ký hiệu đơn lẻ sang chuỗi ký hiệu.


\bibliographystyle{plain}
\bibliography{references}

\appendix
\chapter{Thông số và hướng dẫn tái lập}
Các hyperparameter, phiên bản thư viện, cấu hình phần cứng và lệnh chạy cuối cùng sẽ được cập nhật từ experiment log.\todo{Bổ sung sau khi chốt run thực nghiệm.}

\chapter{Kế hoạch thực hiện}

Thời gian thực hiện dự kiến từ ngày 20/07/2026 đến ngày 23/09/2026, tương đương chín tuần và ba ngày. Kế hoạch dưới đây dùng thuật ngữ thống nhất và giữ tập kiểm thử cho đánh giá cuối cùng.

\begin{table}[ht]
\centering
\small
\caption{Kế hoạch thực hiện dự kiến}
\label{tab:project-plan}
\begin{tabular}{p{2.1cm}p{7.0cm}p{4.0cm}}
\toprule
Thời gian & Công việc & Sản phẩm/tiêu chí liên quan\\
\midrule
20--26/07 & Khảo sát tài liệu; tải, kiểm tra tính toàn vẹn và audit ASL-HG; hoàn thiện đề cương. & Báo cáo audit, citation dataset, đề cương.\\
27/07--09/08 & Xây dựng MediaPipe Hand Landmarker để định vị bàn tay và trích xuất ROI; ghi manifest các trường hợp không phát hiện/ROI không hợp lệ. & Cache ROI và failure manifest.\\
10--16/08 & Chuẩn hóa ảnh; tạo split phân tầng train/validation/test 80/10/10; kiểm tra leakage. & Manifest split tái lập.\\
17--30/08 & Huấn luyện các baseline CNN, MobileNetV4 Conv-S đóng băng và MediaPipe landmarks+SVM; chọn checkpoint bằng validation. & Checkpoint, learning curves, log và artifact tái lập.\\
31/08--06/09 & Fine-tuning toàn bộ MobileNetV4 Conv-S; mọi lựa chọn dùng validation, không dùng test. & Bảng so sánh validation, mô hình ứng viên.\\
07--13/09 & Xây dựng nguyên mẫu desktop OpenCV; đo độ trễ/FPS và xử lý trạng thái không phát hiện bàn tay. & Ứng dụng demo, log latency/FPS.\\
14--20/09 & Chốt mô hình theo validation, đánh giá cuối cùng một lần trên test; lập confusion matrix và phân tích O/0; viết báo cáo. & Metrics test, confusion matrix, bản thảo báo cáo.\\
21--23/09 & Rà soát kỹ thuật, báo cáo, tài liệu tham khảo và chuẩn bị bảo vệ. & Bản nộp, slide và video demo.\\
\bottomrule
\end{tabular}
\end{table}

Mỗi mục tiêu trong Mở đầu được nối với một artifact hoặc chỉ số: pipeline tái lập được kiểm tra bằng manifest/config; mô hình được đánh giá bằng Accuracy, Precision, Recall, F1-score và confusion matrix; O/0 được kiểm tra theo hai hướng nhầm lẫn; nguyên mẫu desktop được đo bằng độ trễ và FPS.


\end{document}
